% ============================================================
% Section 89: 钠金属的费米能
% ============================================================
\section{固体理论：钠金属的费米能}
\label{sec:sodium-fermi-energy}

\begin{modelbox}[题目：钠金属传导电子浓度与费米能的计算]
金属钠是体心立方结构，晶格常数 $a=4.25\times10^{-8}\,\text{cm}$。求传导电子浓度。假定每个原子贡献一个传导电子，用自由电子气模型，推导 $T=0\,$K 时费米能的表达式，并说明它只与传导电子浓度有关而与晶体体积无关。
\end{modelbox}

\begin{tipbox}[考点快速分析]
\begin{itemize}
    \item \textbf{BCC电子浓度}：每个惯用晶胞含 2 个原子，$n=2/a^3$
    \item \textbf{费米能推导}：由态密度积分，总电子数 $N=\int_0^{E_F}g(E)dE$
    \item \textbf{核心结论}：$E_F\propto n^{2/3}$，与 $V$ 无关
    \item \textbf{数值估算}：钠的 $E_F\sim3\,$eV，$T_F\sim3.7\times10^4\,$K
\end{itemize}
\end{tipbox}

\begin{derivationbox}[标准解题过程]
\textbf{Step 1: 传导电子浓度}

钠为体心立方（BCC）结构，每个惯用晶胞含 2 个原子。晶格常数 $a=4.25\times10^{-8}\,\text{cm}=4.25\times10^{-10}\,\text{m}$。

每个钠原子贡献 1 个传导电子，电子数密度
\begin{equation}
n=\frac{2}{a^3}=\frac{2}{(4.25\times10^{-10})^3}
\approx\frac{2}{7.676\times10^{-29}}
\approx\boxed{2.6\times10^{28}\,\text{m}^{-3}.}
\end{equation}

\textbf{Step 2: 绝对零度费米能的推导}

自由电子气模型中，考虑自旋简并，三维态密度
\begin{equation}
g(E)=\frac{V}{2\pi^2}\Bigl(\frac{2m}{\hbar^2}\Bigr)^{3/2}E^{1/2}.
\end{equation}

$T=0$ 时，费米面以下所有态被占据。总电子数
\begin{equation}
N=\int_0^{E_F^0}g(E)\,dE
=\frac{V}{2\pi^2}\Bigl(\frac{2m}{\hbar^2}\Bigr)^{3/2}\int_0^{E_F^0}E^{1/2}\,dE
=\frac{V}{3\pi^2}\Bigl(\frac{2m}{\hbar^2}\Bigr)^{3/2}(E_F^0)^{3/2}.
\end{equation}

解出费米能
\begin{equation}
\boxed{E_F^0=\frac{\hbar^2}{2m}(3\pi^2n)^{2/3}.}
\end{equation}

\textbf{与体积的关系}：$E_F^0$ 仅依赖于电子数密度 $n=N/V$（以及基本常数 $m$ 和 $\hbar$），而与晶体总体积 $V$ 无关。只要 $n$ 不变，无论晶体尺寸如何，费米能均相同。这是自由电子气模型的重要特征，反映了电子的量子简并性。

\textbf{Step 3: 数值计算}

代入 $n=2.6\times10^{28}\,\text{m}^{-3}$：
\begin{align}
E_F^0&=\frac{(1.055\times10^{-34})^2}{2\times9.11\times10^{-31}}\times\bigl(3\pi^2\times2.6\times10^{28}\bigr)^{2/3} \notag\\
&\approx6.11\times10^{-39}\times(7.70\times10^{29})^{2/3} \notag\\
&\approx6.11\times10^{-39}\times8.40\times10^{19} \notag\\
&\approx5.13\times10^{-19}\,\text{J}\approx\boxed{3.2\,\text{eV}.}
\end{align}

对应的费米温度
\begin{equation}
T_F=\frac{E_F^0}{k_B}\approx\frac{3.2\times1.6\times10^{-19}}{1.38\times10^{-23}}\approx3.7\times10^4\,\text{K}.
\end{equation}
远高于室温，因此常温下钠中电子处于高度简并态。
\end{derivationbox}

\begin{formulabox}[答案汇总]
\begin{center}
\begin{tabular}{ll}
\toprule
\textbf{物理量} & \textbf{表达式 / 数值} \\
\midrule
电子浓度 & $n=2/a^3\approx2.6\times10^{28}\,\text{m}^{-3}$ \\
费米能 & $E_F^0=\dfrac{\hbar^2}{2m}(3\pi^2n)^{2/3}\approx3.2\,\text{eV}$ \\
与体积关系 & $E_F^0$ 仅依赖于 $n$，与 $V$ 无关 \\
费米温度 & $T_F\approx3.7\times10^4\,\text{K}$ \\
\bottomrule
\end{tabular}
\end{center}
\end{formulabox}

\begin{insightbox}[物理图像]
\begin{itemize}
    \item 费米能是电子简并压力的体现。即使 $T=0$，电子也因泡利不相容原理而具有巨大的零点动能。金属的刚性、电子的高热导率等性质都与此有关。
    \item $E_F\propto n^{2/3}$ 意味着电子浓度越高，费米能越大。碱金属（如 Na、K）的费米能在 $2\!\sim\!4\,$eV 量级；过渡金属（如 Cu、Ag）因 $d$ 电子贡献，$n$ 更大，$E_F$ 可达 $5\!\sim\!7\,$eV。
    \item ``与体积无关''的结论在热力学极限（$N\to\infty$，$V\to\infty$，$n$ 固定）下严格成立。对于纳米尺度的金属颗粒，量子尺寸效应会导致能级离散化，此时费米能的定义需要修正。
\end{itemize}
\end{insightbox}
